\documentclass[11pt,a4paper]{article}

% --- Encoding and language ---------------------------------------------------
\usepackage[utf8]{inputenc}
\usepackage[T1]{fontenc}
\usepackage{lmodern}

% --- Math and symbols --------------------------------------------------------
\usepackage{amsmath,amssymb,amsthm,amsfonts}
\usepackage{bm}
\usepackage{mathtools}

% --- Layout, figures, tables -------------------------------------------------
\usepackage[margin=1in]{geometry}
\usepackage{graphicx}
\usepackage{float}
\usepackage{caption}
\usepackage{subcaption}
\usepackage{booktabs}
\usepackage{array}
\usepackage{multirow}
\usepackage{enumitem}
\usepackage{microtype}
\usepackage{xcolor}
\usepackage{titlesec}
\usepackage{hyperref}
\hypersetup{
  colorlinks=true,
  linkcolor=blue!50!black,
  urlcolor=blue!60!black,
  citecolor=blue!50!black
}

% --- Caption style -----------------------------------------------------------
\captionsetup{font=small,labelfont=bf,labelsep=period}

% --- Theorem-like environments ----------------------------------------------
\newtheorem{definition}{Definition}
\newtheorem{remark}{Remark}

% --- Custom commands ---------------------------------------------------------
\newcommand{\E}{\mathbb{E}}
\newcommand{\Var}{\mathrm{Var}}
\newcommand{\given}{\,\vert\,}
\newcommand{\dnorm}{\mathcal{N}}
\newcommand{\code}[1]{\texttt{\small #1}}
\newcommand{\figpath}[1]{../outputs/home_advantage/figures/#1}

% --- Title -------------------------------------------------------------------
\title{\textbf{Does the Home Crowd Really Help?}\\[2pt]
\large A Statistical Investigation of the Home-Race Advantage in Formula~1
(1950--2026)}
\author{Applied Statistics Project Report}
\date{\today}

\begin{document}
\maketitle

\begin{abstract}
We investigate whether Formula~1 drivers perform measurably better at their
home Grand Prix than abroad, using the publicly available Ergast/Kaggle F1
dataset spanning 27{,}304 driver--race observations across 1{,}152 races,
865 drivers and 77 circuits from 1950 to 2026. We construct an
\emph{era-adjusted teammate delta} which controls for the dominant confound
in F1 (machinery) and for the historically shifting points system. Three
complementary analyses are reported: (i) a paired $t$-test on per-driver home
vs.\ away means; (ii) a linear mixed-effects model with random
driver-intercepts and a \texttt{home\_race}$\times$\texttt{experience} interaction;
and (iii) an XGBoost regressor evaluated with chronological time-series
cross-validation and explained via SHAP values. The paired test shows
\emph{no significant aggregate home advantage} ($t_{439}=0.65$, $p=0.515$).
The mixed-effects model reveals a more subtle picture: there is a
statistically significant \emph{interaction} between home-race and experience
($\beta_{\text{home}\times\text{exp}}=0.00178$, $p=0.035$): young drivers
underperform at home, while experienced drivers benefit. The XGBoost model
attains a stable normalized RMSE of $0.147\pm 0.009$ across five
chronological folds and SHAP confirms the interaction visually.
\end{abstract}

\tableofcontents
\vspace{1em}

% =============================================================================
\section{Motivation}
% =============================================================================

A central question in sports psychology and sports economics is whether
athletes systematically benefit from competing in front of a friendly home
crowd. In team sports the answer is uncontroversially \emph{yes}: a large body
of work documents home advantages in football, basketball, baseball and
beyond. In Formula~1, however, the answer is far from obvious for two reasons:

\begin{enumerate}[leftmargin=*]
    \item \textbf{The car dominates the driver.} Roughly $70\%$--$80\%$ of
    the variance in finishing position is attributable to the constructor.
    A Mercedes in 2015 finishes ahead of a Marussia regardless of who is
    driving. Any naive home-race analysis is confounded by which constructor
    happens to enter a driver's home race in a given year.

    \item \textbf{The points system has changed many times.} In 1950 a race
    win was worth 9 points; in 1991 it was 10; since 2010 it is 25, with
    additional rules for fastest lap, sprint races, and dropped scores.
    Comparing raw points across eras is meaningless.
\end{enumerate}

A statistically honest analysis must neutralise both confounds. We do so via
an \emph{era-adjusted teammate delta} (Section~\ref{sec:target}), then layer
three increasingly nuanced analyses on top of it.

\paragraph{Research questions.}
\begin{description}[leftmargin=*,style=nextline]
    \item[RQ1] Do drivers, on average, perform better at home than away?
    \item[RQ2] Does any home effect depend on driver characteristics --
    in particular, on \emph{experience}?
    \item[RQ3] Can we build a predictive model whose out-of-time accuracy
    and feature attributions corroborate the inferential findings?
\end{description}

% =============================================================================
\section{Data}\label{sec:data}
% =============================================================================

We use the well-known Ergast Formula~1 database
(re-distributed on Kaggle), comprising 14 relational CSVs. After joining
\code{results}, \code{races}, \code{drivers}, \code{circuits} and
\code{status}, and after enriching driver nationality with a custom
\code{nationality $\to$ country} mapping, the working frame contains:

\begin{center}
\begin{tabular}{lr}
\toprule
Quantity & Value \\
\midrule
Driver--race observations & 27{,}304 \\
Distinct drivers          & 865 \\
Distinct races            & 1{,}152 \\
Distinct circuits         & 77 \\
Years covered             & 1950--2026 \\
Home starts ($\code{home\_race}=1$)  & 2{,}545 \\
Away starts ($\code{home\_race}=0$)  & 24{,}759 \\
\bottomrule
\end{tabular}
\end{center}

The strong class imbalance ($\approx 9.3\%$ home starts) reflects a structural
feature of F1: most drivers are British, German or Italian, but most circuits
sit elsewhere on the calendar in a given season.

\subsection{Defining ``home race''}
A driver's nationality is recorded as a free-text label (e.g.\ \textit{British},
\textit{Argentine-Italian}, \textit{Monegasque}), while circuit countries
follow ISO-style short names (\textit{UK}, \textit{Argentina}, \textit{Monaco}).
We resolve this mismatch with a curated mapping
$\phi:\text{Nationality}\to\text{Country}$ shipped as \code{countries.csv}.
The home-race indicator is then:
\begin{equation}
    \texttt{home\_race}_{i,r}
    \;=\;
    \mathbf{1}\!\big[\,\phi(\texttt{nationality}_i)
    \;=\;\texttt{country}(\texttt{circuit}_r)\,\big].
\end{equation}
We normalise the U.S.\ circuit-country variants (\textit{``United States''}
$\mapsto$ \textit{``USA''}) so that an American driver is correctly tagged
as home for both Indianapolis and Austin events.

\subsection{Sample-size sanity check}
Figure~\ref{fig:nationalities} shows the top twelve nationalities by distinct
drivers. British and Italian drivers dominate; this matters because they also
have the most home opportunities and therefore the most statistical power in
any per-driver paired test.

\begin{figure}[H]
    \centering
    \includegraphics[width=0.78\textwidth]{\figpath{04_top_nationalities.png}}
    \caption{Top twelve nationalities by distinct drivers. The British and
    Italian populations dominate the home-race observations.}
    \label{fig:nationalities}
\end{figure}

% =============================================================================
\section{Target Variable: Era-Adjusted Teammate Delta}\label{sec:target}
% =============================================================================

\subsection{Why teammate-relative performance?}
The single largest source of variance in a race result is \emph{the car}.
Two drivers in the same constructor share, up to mechanical luck, the same
machinery. Their \emph{within-team} difference therefore isolates driver-level
skill from car-level dominance. This is a classical design choice in F1
analytics and serves as our identification strategy: by comparing each
driver to their own teammate in the same race, the constructor effect
cancels out.

\subsection{Why normalise by the race's maximum?}
Raw points are not comparable across eras. A 1950 winner scored 9 points;
a 2024 winner scores 25 plus 1 for fastest lap. We therefore divide each
driver's points by the maximum scored at that specific race:
\begin{equation}
    p^{\star}_{i,r} \;=\; \frac{p_{i,r}}{\max_{j\in r} p_{j,r}}
    \;\in\;[0,1].
\end{equation}
This converts every race onto a common $[0,1]$ scale -- $1.0$ means the
driver won, $0.5$ means they scored half what the winner scored, $0$ means
no points -- regardless of era.

\subsection{Construction of the teammate delta}
For driver $i$ in race $r$, let $T(i,r)$ be the set of drivers fielded by
$i$'s constructor in race $r$ (the team mean is taken over both drivers,
which simplifies algebra; the difference from a true ``other-teammate'' delta
is constant in expectation and absorbed by the intercept). Then
\begin{equation}\label{eq:delta}
    \Delta_{i,r}
    \;=\;
    p^{\star}_{i,r}
    \;-\;
    \frac{1}{|T(i,r)|}\sum_{j\in T(i,r)} p^{\star}_{j,r}.
\end{equation}
Positive $\Delta_{i,r}$ means the driver outperformed their team's average
for that race; negative means underperformed.

\begin{remark}
Because the team mean is computed over both drivers, $\Delta_{i,r}$ for a
two-car team is exactly half the gap to the teammate. The scale changes but
the sign and the statistical conclusions do not.
\end{remark}

\subsection{Distribution of the target}
Figure~\ref{fig:delta-hist} shows the empirical distribution of
$\Delta_{i,r}$. It is symmetric and tightly concentrated near zero, as
expected: the average driver, by construction, performs at the team's
average. The long tails reflect blow-out races where one car dominates the
other (e.g.\ DNFs that score zero against a winning teammate).

\begin{figure}[H]
    \centering
    \includegraphics[width=0.78\textwidth]{\figpath{01_teammate_delta_hist.png}}
    \caption{Empirical distribution of the era-adjusted teammate delta
    $\Delta_{i,r}$. The mean is approximately zero by construction;
    deviations from zero quantify driver-specific over- or under-performance.}
    \label{fig:delta-hist}
\end{figure}

% =============================================================================
\section{Auxiliary Features}
% =============================================================================

Beyond the target, we engineer three covariates used in the mixed-effects
model and the XGBoost predictor.

\paragraph{Experience.}
For each driver $i$, let $y_i^{\min}$ be the year of their first recorded
race. Then $\texttt{experience\_years}_{i,r} = \texttt{year}_r - y_i^{\min}$.
This is the standard ``years since debut'' measure and ranges from $0$ for
a rookie to upward of $20$ for long-running careers.

\paragraph{Grid position.}
Qualifying position is the strongest non-car predictor of race result;
including it controls for one-off mechanical/qualifying problems that have
nothing to do with the home crowd.

\paragraph{Baseline-normalised crash risk.}
Some circuits are simply more dangerous than others (Monaco walls, Spa weather).
For each driver--race we compute
\begin{equation}
    \texttt{is\_crash}_{i,r} = \mathbf{1}\!\big[\texttt{statusId}_{i,r}\in\{3,4\}\big],
\end{equation}
i.e.\ accident or collision. The circuit-level baseline is
\begin{equation}
    c_r = \frac{1}{n_r}\sum_{i\in r} \texttt{is\_crash}_{i,r},
\end{equation}
and the residual is $\texttt{normalised\_crash\_risk}_{i,r} =
\texttt{is\_crash}_{i,r} - c_{\text{circuit}(r)}$.

Figure~\ref{fig:crash} shows the empirical circuit crash rates (circuits
with $\ge 100$ starts), confirming that the worst tracks have $\sim 30\%$
of starters retire in incidents -- a baseline that absolutely must be
controlled for.

\begin{figure}[H]
    \centering
    \includegraphics[width=0.88\textwidth]{\figpath{05_circuit_crash_rates.png}}
    \caption{Top fifteen most dangerous circuits (accident or collision per
    start, minimum 100 starts). Circuit-baseline crash rates exceed $25\%$
    for the worst tracks.}
    \label{fig:crash}
\end{figure}

% =============================================================================
\section{Inferential Analysis I: Paired \texorpdfstring{$t$}{t}-Test}
% =============================================================================

\subsection{Why \emph{paired}?}
An unpaired (Welch) comparison of \emph{all} home rows against \emph{all}
away rows would conflate two effects:
\begin{enumerate}[label=(\roman*),leftmargin=*]
    \item the genuine home-vs-away difference within drivers, and
    \item the fact that some drivers are simply better than others and just
    happen to race at home more often (a between-driver confounder).
\end{enumerate}
By computing \emph{per-driver} means and pairing them, we eliminate (ii).
Each driver acts as their own control.

\subsection{Statistical setup}
For every driver $i$ who ever started a home race we form
\begin{equation}
\bar{\Delta}_i^{\text{home}} = \frac{1}{|H_i|}\sum_{r\in H_i}\Delta_{i,r},
\qquad
\bar{\Delta}_i^{\text{away}} = \frac{1}{|A_i|}\sum_{r\in A_i}\Delta_{i,r},
\end{equation}
where $H_i,A_i$ are driver $i$'s home and away races. The driver-level
difference is $d_i = \bar{\Delta}_i^{\text{home}} - \bar{\Delta}_i^{\text{away}}$.
Under the null hypothesis $H_0:\E[d_i]=0$ (no home advantage), the test
statistic
\begin{equation}
    t \;=\; \frac{\bar{d}}{s_d/\sqrt{n}} \;\sim\; t_{n-1}
\end{equation}
follows Student's $t$ distribution with $n-1$ degrees of freedom, where
$\bar{d}$ and $s_d$ are the sample mean and standard deviation of the
$\{d_i\}$ and $n$ is the number of drivers with at least one home
\emph{and} one away race.

\subsection{Results}
\begin{center}
\begin{tabular}{lr}
\toprule
Quantity & Value \\
\midrule
Number of paired drivers ($n$)             & $440$ \\
Mean paired difference ($\bar{d}$)         & $+0.00265$ \\
SD of paired differences ($s_d$)           & $0.08550$ \\
$95\%$ CI for $\bar{d}$                    & $[-0.00534,\; +0.01064]$ \\
$t$-statistic                              & $0.651$ \\
Degrees of freedom                         & $439$ \\
$p$-value (two-sided)                      & $0.515$ \\
Cohen's $d$                                & $0.031$ \\
\bottomrule
\end{tabular}
\end{center}

The $95\%$ confidence interval narrowly straddles zero and the effect size
is negligible. We \emph{do not reject} the null: at the aggregate level there
is no detectable home advantage in F1.

\subsection{Visual diagnostics}
The per-driver effects (Fig.~\ref{fig:driver-effect}) and the paired-difference
histogram (Fig.~\ref{fig:diff-hist}) make the conclusion concrete: while
some drivers do measurably better at home and others measurably worse, the
distribution is symmetric around zero.

\begin{figure}[H]
    \centering
    \includegraphics[width=0.72\textwidth]{\figpath{06_per_driver_effect.png}}
    \caption{Per-driver home-minus-away mean teammate delta, top 25 drivers
    by total starts. Red bars are positive (driver helped by home crowd),
    blue bars are negative (driver hurt by home crowd).}
    \label{fig:driver-effect}
\end{figure}

\begin{figure}[H]
    \centering
    \includegraphics[width=0.78\textwidth]{\figpath{07_paired_diff_hist.png}}
    \caption{Histogram of per-driver paired differences $d_i$. The red
    vertical line marks the mean ($+0.0027$). The distribution is
    near-symmetric about zero -- a textbook ``no effect'' picture.}
    \label{fig:diff-hist}
\end{figure}

\paragraph{Comparison with raw distributions.} Figures~\ref{fig:box} and
\ref{fig:meanci} present the same comparison without pairing. They are
visually identical between groups and underscore the same message: no
aggregate effect.

\begin{figure}[H]
    \centering
    \begin{subfigure}{0.48\textwidth}
        \centering
        \includegraphics[width=\linewidth]{\figpath{02_home_vs_away_box.png}}
        \caption{Boxplots.}
        \label{fig:box}
    \end{subfigure}
    \hfill
    \begin{subfigure}{0.48\textwidth}
        \centering
        \includegraphics[width=\linewidth]{\figpath{03_home_vs_away_meanci.png}}
        \caption{Means with $95\%$ CIs.}
        \label{fig:meanci}
    \end{subfigure}
    \caption{Unpaired comparison of $\Delta_{i,r}$ across home and away.}
    \label{fig:hvsa}
\end{figure}

% =============================================================================
\section{Inferential Analysis II: Linear Mixed-Effects Model}
% =============================================================================

\subsection{Why not stop at the $t$-test?}
The paired $t$-test discards information by collapsing each driver to two
numbers. It also implicitly assumes the home effect is constant across
drivers. A more powerful approach is a \emph{linear mixed-effects model}
(LMM, also known as GLMM with identity link and Gaussian noise) that

\begin{itemize}[leftmargin=*]
    \item models every individual driver--race observation,
    \item lets each driver have their own baseline via a random intercept,
    \item allows the home-race effect to \emph{depend on driver experience}.
\end{itemize}

\subsection{Model specification}
For driver $i$ in race $r$ we posit
\begin{align}
    \Delta_{i,r}
    \;=\;&
    \beta_0
    \;+\;\beta_1 H_{i,r}
    \;+\;\beta_2 E_{i,r}
    \;+\;\beta_3 (H_{i,r}\!\cdot\! E_{i,r}) \nonumber\\[2pt]
    & + u_i + \varepsilon_{i,r},
    \label{eq:lmm}
\end{align}
where $H_{i,r}=\texttt{home\_race}_{i,r}$, $E_{i,r}=\texttt{experience\_years}_{i,r}$,
the random intercept satisfies $u_i\sim\dnorm(0,\sigma_u^2)$, and the residual
$\varepsilon_{i,r}\sim\dnorm(0,\sigma^2)$ is assumed independent of
$u_i$. Estimation is by restricted maximum likelihood (REML) via
\code{statsmodels.MixedLM}.

\paragraph{Interpretation of coefficients.}
\begin{itemize}[leftmargin=*]
    \item $\beta_0$: expected teammate delta of a rookie ($E=0$) racing away.
    \item $\beta_1$: home-race premium for a rookie ($E=0$).
    \item $\beta_2$: change in teammate delta per extra year of experience,
    \emph{when racing away}.
    \item $\beta_3$: \emph{differential} slope of experience for home races.
    The total experience slope at home is $\beta_2+\beta_3$.
    \item $\sigma_u^2$: between-driver variance left unexplained by fixed
    effects -- the ``how much do drivers genuinely differ'' term.
    \item $\sigma^2$: within-driver residual variance.
\end{itemize}

\subsection{Estimation results}

\begin{center}
\renewcommand{\arraystretch}{1.15}
\begin{tabular}{lrrrr}
\toprule
Term & Estimate & Std.\ Error & $z$ & $p$-value \\
\midrule
Intercept $\beta_0$                     & $-0.00984$ & $0.00249$ & $-3.96$ & $7.6\times 10^{-5}$ \\
\code{home\_race} $\beta_1$             & $-0.00613$ & $0.00466$ & $-1.32$ & $0.188$ \\
\code{experience\_years} $\beta_2$      & $-0.00143$ & $0.00029$ & $-4.96$ & $7.2\times 10^{-7}$ \\
\code{home\_race$\times$experience} $\beta_3$
                                        & $+0.00178$ & $0.00084$ & $+2.10$ & $0.035$ \\
\midrule
Group (driver) variance $\sigma_u^2$    & $0.00174$ & --- & --- & --- \\
Residual variance $\sigma^2$            & $0.02279$ & --- & --- & --- \\
$N$ observations                        & $27{,}304$ & --- & --- & --- \\
$N$ groups (drivers)                    & $865$      & --- & --- & --- \\
\bottomrule
\end{tabular}
\end{center}

\subsection{What the model says}

\begin{enumerate}[leftmargin=*]
    \item \textbf{The constant home main-effect $\beta_1$ is not significant}
    ($p=0.188$). This is consistent with the paired $t$-test.
    \item \textbf{Experience has a small, highly significant \emph{negative}
    main slope away from home} ($\beta_2=-0.00143$, $p<10^{-6}$). The most
    plausible reading is statistical, not biological: by construction
    $\Delta_{i,r}$ is centred at zero, so as drivers accumulate experience
    they often face increasingly young teammates whose averages they help
    define. The size of the effect is tiny.
    \item \textbf{The home$\times$experience interaction is positive and
    significant} ($\beta_3=+0.00178$, $p=0.035$). The total experience
    slope \emph{at home} is $\beta_2+\beta_3=+0.00035$, i.e.\ effectively
    flat -- whereas at away races experienced drivers slowly lose ground
    to younger teammates. Equivalently, the home crowd appears to
    \emph{neutralise} the age-related drift.
\end{enumerate}

\paragraph{Intraclass correlation.}
The share of variance attributable to driver identity is
\begin{equation}
\rho = \frac{\sigma_u^2}{\sigma_u^2+\sigma^2}
     = \frac{0.00174}{0.00174+0.02279}
     \approx 0.071,
\end{equation}
i.e.\ roughly $7\%$ of residual variance is between-driver. This is small
but non-negligible and justifies the use of a mixed model: ordinary least
squares with naive standard errors would understate the uncertainty on
$\beta_1$.

\subsection{Visualising the interaction}
Figure~\ref{fig:lmm-marginal} plots the model-predicted teammate delta as a
function of experience, separately at home and away. The two lines visibly
cross: rookies are slightly \emph{worse} at home (perhaps too eager,
unfamiliar with pressure), while veterans are slightly \emph{better} at home
relative to away.

\begin{figure}[H]
    \centering
    \includegraphics[width=0.8\textwidth]{\figpath{08_mixedlm_marginal.png}}
    \caption{Predicted teammate delta from the mixed-effects model
    (Eq.~\eqref{eq:lmm}) as a function of experience, holding all other
    quantities at their mean. The red ``home'' line is essentially flat;
    the blue ``away'' line slopes gently downward.}
    \label{fig:lmm-marginal}
\end{figure}

% =============================================================================
\section{Predictive Analysis: XGBoost + SHAP + Time-Series Backtest}
% =============================================================================

The inferential models test \emph{structured hypotheses}. We additionally
fit a flexible non-parametric predictor to (i) test how much variance can
be explained out-of-sample at all, and (ii) discover non-linear interactions
that a linear model cannot see.

\subsection{Model}
We train an XGBoost regressor:
\begin{equation}
    \hat{\Delta}_{i,r}
    \;=\;
    f_{\text{XGB}}(E_{i,r},\,H_{i,r},\,G_{i,r},\,c_{r}),
\end{equation}
with $G_{i,r}$ the grid position and $c_r$ the circuit baseline crash rate.
Hyperparameters: $100$ trees, max depth $4$, learning rate $0.1$, fixed seed.

\subsection{Why time-series cross-validation?}
F1 is a chronological process: cars, regulations and driver populations
drift. Random $k$-fold cross-validation would let the model peek into the
future of the same driver/season -- a leakage that systematically overstates
performance. We instead use \code{TimeSeriesSplit} with five expanding
windows. In every fold, training data strictly precedes test data in
absolute time. Concretely:

\begin{center}
\begin{tabular}{rllrrrr}
\toprule
Fold & Train years & Test years & $n_{\text{train}}$ & $n_{\text{test}}$ & RMSE & MAE \\
\midrule
1 & 1950--1971 & 1971--1982 & 4{,}549  & 4{,}547 & 0.162 & 0.106 \\
2 & 1950--1982 & 1982--1992 & 9{,}096  & 4{,}547 & 0.137 & 0.081 \\
3 & 1950--1992 & 1992--2004 & 13{,}643 & 4{,}547 & 0.145 & 0.086 \\
4 & 1950--2004 & 2004--2015 & 18{,}190 & 4{,}547 & 0.152 & 0.099 \\
5 & 1950--2015 & 2015--2026 & 22{,}737 & 4{,}547 & 0.140 & 0.093 \\
\midrule
\textbf{Mean} & & & & & \textbf{0.147} & \textbf{0.093} \\
\textbf{SD}   & & & & & 0.009 & 0.009 \\
\bottomrule
\end{tabular}
\end{center}

\begin{figure}[H]
    \centering
    \includegraphics[width=0.82\textwidth]{\figpath{13_backtest_folds.png}}
    \caption{Time-series backtest: RMSE and MAE per chronological fold.
    The model generalises remarkably stably across eras.}
    \label{fig:backtest}
\end{figure}

\paragraph{How good is this in real F1 terms?}
The error is in the normalised $[0,1]$ scale. Multiplying by the modern
race maximum of $25$ points, a normalised MAE of $0.093$ corresponds to
roughly $\pm 2.3$ raw points in expectation -- about the difference
between finishing $9^{\text{th}}$ and $7^{\text{th}}$ in a contemporary
points system. Given that we use only four features, this is a strong
signal.

\subsection{Feature importance}
\begin{figure}[H]
    \centering
    \includegraphics[width=0.72\textwidth]{\figpath{09_xgb_importance.png}}
    \caption{XGBoost gain-importance of the four predictors. Grid position
    dominates by a wide margin, as expected; experience and circuit crash
    rate make secondary contributions; \texttt{home\_race} is the weakest.}
    \label{fig:imp}
\end{figure}

\subsection{SHAP analysis}
SHAP (SHapley Additive exPlanations) decomposes each prediction into a sum
of additive contributions from each feature. The beeswarm in
Fig.~\ref{fig:shap-beeswarm} shows, for a random subsample of 5{,}000 rows,
the magnitude and direction of each feature's contribution.

\begin{figure}[H]
    \centering
    \includegraphics[width=0.78\textwidth]{\figpath{10_shap_beeswarm.png}}
    \caption{SHAP beeswarm. Each dot is one driver--race; horizontal position
    is that feature's contribution to the prediction; colour encodes the
    feature value. Grid is the dominant driver and behaves monotonically.}
    \label{fig:shap-beeswarm}
\end{figure}

The dependence plot for \texttt{experience\_years}, coloured by
\texttt{home\_race} (Fig.~\ref{fig:shap-dep-exp}), corroborates the linear
mixed-effects finding: at any fixed experience level the home dots (red)
sit slightly higher than the away dots (blue) for experienced drivers,
matching the positive interaction $\beta_3$.

\begin{figure}[H]
    \centering
    \includegraphics[width=0.78\textwidth]{\figpath{11_shap_dependence_exp_home.png}}
    \caption{SHAP dependence: SHAP value of \texttt{experience\_years} vs.\
    experience, coloured by \texttt{home\_race}. Red = home, blue = away.
    The vertical separation between red and blue grows with experience --
    visual confirmation of the positive interaction term.}
    \label{fig:shap-dep-exp}
\end{figure}

\begin{figure}[H]
    \centering
    \includegraphics[width=0.78\textwidth]{\figpath{12_shap_dependence_grid.png}}
    \caption{SHAP dependence on grid position. The model has learned the
    obvious monotone relationship -- worse starting position predicts
    larger negative teammate deltas -- but with a clear non-linear curve
    in the front of the grid.}
    \label{fig:shap-dep-grid}
\end{figure}

% =============================================================================
\section{Discussion}
% =============================================================================

\subsection{Convergent evidence across three methods}
All three analyses tell the same story from different angles:
\begin{itemize}[leftmargin=*]
    \item \textbf{Paired $t$-test}: no aggregate home effect on driver
    performance.
    \item \textbf{Mixed-effects model}: no main effect of home; \emph{but} a
    significant interaction with experience that flattens the
    age-related drift.
    \item \textbf{XGBoost + SHAP}: \texttt{home\_race} carries little raw
    importance, but its SHAP-level interaction with experience is visible
    in the dependence plot, mirroring $\beta_3$ from the LMM.
\end{itemize}

\subsection{Why this is interesting}
The conventional wisdom -- that home crowds confer a measurable boost --
is \emph{not supported} in Formula~1 at the aggregate level. The much-
hyped emotional advantage of a driver's home Grand Prix appears, in the
cold light of 76 seasons of data, to be roughly zero \emph{on average}.

But the picture is more nuanced than a flat null. Young drivers seem to
do slightly \emph{worse} at home, perhaps consistent with media pressure
and family-and-friends distraction. Veterans seem to do slightly
\emph{better} -- which we interpret cautiously as a small experience-
moderated effect rather than a strong causal claim, given $p=0.035$ on a
secondary interaction term.

\subsection{Limitations}
\begin{enumerate}[leftmargin=*]
    \item \textbf{Class imbalance.} Only $9.3\%$ of starts are home starts;
    statistical power on $\beta_1$ is therefore modest.
    \item \textbf{Causal language.} Observational data cannot fully resolve
    confounders such as constructor selection (a team may field a stronger
    car at a home race for sponsorship reasons).
    \item \textbf{Heterogeneous nationalities.} Some labels
    (\textit{Argentine-Italian}, \textit{Rhodesian}, \textit{East German})
    require mapping decisions; sensitivity analyses are left for future
    work.
    \item \textbf{Static features.} We do not include weather, tyre, or
    in-season form -- all of which a more granular study could add.
\end{enumerate}

\subsection{Take-home}
A statistically careful analysis of 76 years of F1 data finds
\emph{no detectable average home advantage}, with a small but real
experience-moderated component. The methodological lesson is broader:
the construction of an \emph{era-adjusted, teammate-relative} target is
what allows three otherwise-incomparable statistical techniques (a
classical paired test, a hierarchical linear model and a non-parametric
boosted-tree predictor) to converge on the same substantive answer.

% =============================================================================
\appendix
\section{Reproducibility}
% =============================================================================

The analysis is driven by a single entry point:
\begin{verbatim}
python run_analysis.py
\end{verbatim}
which sequentially executes the modules under \code{src/}:
\code{data\_loader.py} (joining the 14 CSVs), \code{features.py} (engineering
the target and covariates), \code{statistics.py} (paired test and mixed
model), and \code{models.py} (XGBoost + backtest + SHAP). Figures are
written to \code{outputs/figures/}, tables to \code{outputs/tables/}, and a
machine-readable summary to \code{outputs/tables/results.json}.

\paragraph{Software.}
Python 3.12 with \code{pandas}, \code{scipy}, \code{statsmodels},
\code{xgboost} 3.x, \code{shap}, \code{scikit-learn}, \code{matplotlib},
\code{seaborn}. Random seeds fixed to $42$ for XGBoost.

\paragraph{Data integrity.}
The nationality $\to$ country mapping in \code{dataset/countries.csv} is
one-to-one (no row duplication on merge). U.S.\ circuit countries are
normalised to a single \textit{USA} label inside \code{data\_loader.py}.

\end{document}
